\documentclass[12pt]{article}
\usepackage[margin=1in]{geometry}
\usepackage{titlesec}
\usepackage{hyperref}
\usepackage{booktabs}
\usepackage{array}
\usepackage{longtable}
\usepackage{parskip}
 
\title{Software Design Document\\
\large Workflow Queue (WFQ) Administrator Application}
\author{Ivonne Arizpe \and Hector Corona \and Ruby Quinones \and Janise Williams\\[0.5em]
\textit{CS3338 Course Group 9}}
\date{Version 1.0 \\ May 14, 2026}
 
\begin{document}
 
\maketitle
\newpage
\tableofcontents
\newpage
 
% ─────────────────────────────────────────────
\section*{Revision History}
% ─────────────────────────────────────────────
\begin{center}
\begin{tabular}{llll}
\toprule
Name & Date & Reason For Changes & Version \\
\midrule
Ivonne Arizpe & 5/2/2026 & Initial Document & 1.0 \\
Ruby Quinones & 5/2/2026 & Section 3 & 1.2 \\
Ivonne Arizpe & 5/3/2026 & Section 1 Introduction, Section 4 System Architecture & 1.2 \\
Hector Corona & 5/3/2026 & Section 8: Database Design completion & 1.3 \\
Ivonne Arizpe & 5/3/2026 & Section 4, 6, and 7, 9 & 1.4 \\
Janise Williams & 5/4/2026 & Section 9 User Interface & 1.5 \\
Ruby Quinones & 5/6/2026 & Section 10: Requirements Validation and Table of Contents & 1.5 \\
Ivonne Arizpe & 5/6/2026 & Section 12: References & 1.5 \\
Janise Williams & 5/7/2026 & Section 5: Policies and Tactics, Section 11: Glossary & 1.5 \\
Hector Corona & 5/8/2026  & Finalized Section 2, Removed API info & 1.6 \\
Hector Corona & 5/8/2026  & Section 1 and Section 4 Revisions & 1.6 \\
Ruby Quinones & 5/9/2026  & Section 12.2 Regulatory References Revision & 1.7 \\
Ivonne Arizpe & 5/9/2026  & Section 12.2 Regulatory Revisions & 1.7 \\
Janise Williams & 5/9/2026  & Section 10 Requirements Validation Revisions & 1.7 \\
Hector Corona & 5/10/2026  & Added existing architecture info & 1.8 \\
Ruby Quinones & 5/10/2026 & Removed Section 12.2, updated Section 2.2 and 3.4  & 1.8 \\
Janise Williams & 5/10/2026 & Section 10 edited, removed API calling & 1.8 \\
Ivonne Arizpe & 5/10/2026 & Section 1 Introduction updated component descriptions & 1.8 \\
Ruby Quinones & 5/10/2026 & Removed Business Logic Mentions in 6.2 & 1.8 \\
Janise Williams & 5/11/2026 & Revised section 6.4.6 and Section 4.1 DFD & 1.8 \\
Janise Williams & 5/12/2026 & Revised formatting for sections 7 and 8 & 1.8 \\
\bottomrule
\end{tabular}
\end{center}
\newpage
 
% ─────────────────────────────────────────────
\section{Introduction}
% ─────────────────────────────────────────────
 
\subsection{Purpose}
This Software Design Document (SDD) describes the design decisions made by CS3338 Group 9
for the Workflow Queue (WFQ) Administrator Application. It explains how our team structured
and planned the project, based on the ten tasks tracked in our Jira project board
(FinalProject-Scrum). This document is intended to show how our design choices address the
goals of each Jira task.
 
\subsection{Scope}
The WFQ Administrator Application is a web-based tool that allows authorized administrators
to manage workflows, steps, and tasks. Our team's design covers the following areas, each
derived directly from our Jira backlog:
 
\begin{itemize}
    \item Database connectivity and error handling (FS-1, FS-9)
    \item Repository and version control setup (FS-2, FS-6)
    \item LaTeX documentation (FS-3, FS-4)
    \item Docker-based deployment (FS-5)
    \item TestRail integration for testing (FS-7)
    \item User authentication (FS-8)
    \item Bug tracking for login issues (FS-10)
\end{itemize}
 
\subsection{Intended Audience}
\begin{itemize}
    \item \textbf{Group 9 Developers:} For implementation reference.
    \item \textbf{Course Instructor:} To evaluate the team's design decisions and task
    planning.
    \item \textbf{QA Testers:} To understand the system design when writing and executing
    test cases.
\end{itemize}
 
% ─────────────────────────────────────────────
\section{System Overview}
% ─────────────────────────────────────────────
 
The WFQ Administrator Application was planned and tracked using a Jira Scrum board with ten
tasks (FS-1 through FS-10). Each task represents a component or area of work that our team
designed and is responsible for. The system is built on ASP.NET Core MVC with a SQL Server
database backend, containerized using Docker, and tested through TestRail.
 
The design of the system follows a layered approach:
\begin{itemize}
    \item \textbf{Presentation Layer:} The web-based user interface served by ASP.NET Core MVC.
    \item \textbf{Business Logic Layer:} C\# controllers that handle application logic and
    route requests.
    \item \textbf{Data Access Layer:} Entity Framework Core communicating with the SQL Server
    database.
    \item \textbf{Infrastructure Layer:} Docker containers for the frontend, backend, and
    database services.
\end{itemize}
 
% ─────────────────────────────────────────────
\section{Design Decisions by Jira Task}
% ─────────────────────────────────────────────
 
This section describes the design decision and approach for each task on our Jira board.
 
\subsection{FS-1: Fix Database Connection}
The database connection is configured through the application's \texttt{appsettings.json}
file using a named connection string. The connection string points to the SQL Server
container defined in our \texttt{docker-compose.yml}. Entity Framework Core is used as the
ORM to manage all database interactions.
 
To make the connection reliable, the backend service in Docker Compose uses a
\texttt{depends\_on} directive to ensure the database container is running before the
application starts. Connection credentials are stored in environment variables and are not
hardcoded in the source code.
 
\subsection{FS-2: Repository Setup}
The GitHub repository is organized with the following folder structure:
 
\begin{itemize}
    \item \texttt{/docs} --- Contains all \LaTeX{} documentation files including the SRS,
    SDD, Design Specification, Snapshot Objectives, and User Manual.
    \item \texttt{/testrail} --- Contains TestRail test run reports for Snapshots 2, 3,
    and 4 in Markdown format.
    \item \texttt{README.md} --- The project's user manual providing an overview, Jira link,
    and access instructions.
    \item \texttt{docker-compose.yml} --- The Docker Compose configuration for running
    all services.
\end{itemize}
 
All team members were given collaborator access to the repository and committed work
throughout the project lifecycle.
 
\subsection{FS-3: LaTeX SRS (Version 1)}
The Software Requirements Specification was written in \LaTeX{} and stored at
\texttt{/docs/SRS/SRS.tex}. It defines the functional and non-functional requirements for
the WFQ application from the perspective of our team's planned work. The SRS was written
by our team and is based on the expected behaviors described in our TestRail test cases.
 
\subsection{FS-4: LaTeX SDD (Version 2)}
This document (the SDD) is the deliverable for FS-4. It is written in \LaTeX{} and stored
at \texttt{/docs/SSD/SDD.tex}. It describes how the system was designed and how each Jira
task maps to a design decision. The SDD is written entirely by our team and reflects our
own project planning and design choices.
 
\subsection{FS-5: Establish Docker}
The Docker setup is defined in \texttt{docker-compose.yml} at the root of the repository.
The configuration includes three services:
 
\begin{itemize}
    \item \textbf{wfq-frontend} --- A Node.js container serving the frontend of the
    application on port 3000.
    \item \textbf{wfq-admin-app} --- An ASP.NET Core container running the backend
    application on port 8080.
    \item \textbf{wfq-database} --- A Microsoft SQL Server 2019 container running on
    port 1433, with a persistent volume for data storage.
\end{itemize}
 
The frontend depends on the backend, and the backend depends on the database, ensuring
services start in the correct order. Database credentials are passed through environment
variables.
 
\subsection{FS-6: Start Git Repo}
The GitHub repository was initialized at the start of the project with an initial commit
containing the folder structure and README. The main branch serves as the primary branch
for all finalized work. Team members worked on updates and committed directly to main for
finalized changes. The repository is hosted at:
\url{https://github.com}
 
\subsection{FS-7: TestRail Connection}
TestRail was used to manage and execute all test cases for the project. Test runs were
created for Snapshots 2, 3, and 4. After each test run, a summary report was exported and
saved in the \texttt{/testrail} directory of the repository. The TestRail reports include
test case titles, expected results, and Pass/Fail status for each test.
 
\subsection{FS-8: Authentication}
User authentication is handled by the application before granting access to any
administrative page. The design uses a URL query string parameter
(\texttt{?is\_admin=true}) to determine admin authorization, which is checked on every
request. If a user is not authorized, they are redirected to an Access Denied page. This
approach was chosen for simplicity within the scope of the course project.
 
\subsection{FS-9: SQL Error}
SQL error handling is implemented at the data access layer. All database calls are wrapped
in try-catch blocks that catch \texttt{SqlException} and general \texttt{Exception} types.
When an error is caught, the system logs the details to the application log and displays a
user-friendly error message on screen. The application does not crash or expose raw SQL
error details to the end user.
 
\subsection{FS-10: Potential Bug --- User Login Issues with Credentials}
This Jira item was created as a bug ticket to track a known area of weakness: user login
may fail when certain types of credentials are entered, such as those containing special
characters or empty fields. This bug was logged with Highest priority because login is a
critical entry point for the system.
 
The design response to this bug is to ensure that the authentication layer validates and
sanitizes all credential inputs before processing. Error messages shown to the user must be
informative but must not reveal system-level details.
 
% ─────────────────────────────────────────────
\section{Technology Stack}
% ─────────────────────────────────────────────
 
The following technologies were selected by our team for this project:
 
\begin{center}
\begin{tabular}{ll}
\toprule
Technology & Purpose \\
\midrule
ASP.NET Core MVC & Web application framework for the backend and UI \\
C\# & Primary programming language \\
Entity Framework & ORM for database access \\
SQL Server 2019 & Relational database for data storage \\
Docker & Containerization of all services \\
Docker Compose & Orchestration of multi-container deployment \\
GitHub & Version control and collaboration \\
Jira & Task and bug tracking \\
TestRail & Test case management and execution \\
\LaTeX{} & Documentation formatting \\
\bottomrule
\end{tabular}
\end{center}
 
% ─────────────────────────────────────────────
\section{Constraints and Limitations}
% ─────────────────────────────────────────────
 
\begin{itemize}
    \item The application is a course project and is not deployed to a live production
    environment.
    \item Authentication is simplified using a URL query string flag rather than a full
    identity provider integration.
    \item The database schema used is based on an existing project template provided as
    part of the course.
    \item The project scope is limited to what is trackable within the six-snapshot course
    timeline.
\end{itemize}
 
% ─────────────────────────────────────────────
\appendix
\section{Jira Task Summary}
% ─────────────────────────────────────────────
 
\begin{center}
\begin{tabular}{lll}
\toprule
Jira ID & Task & Assignee \\
\midrule
FS-1 & Fix Database Connection & Hector Corona \\
FS-2 & Repository Setup & Ruby Quinones \\
FS-3 & LaTeX SRS (Version 1) & Ivonne Arizpe \\
FS-4 & LaTeX SDD (Version 2) & Janise Williams \\
FS-5 & Establish Docker & Janise Williams \\
FS-6 & Start Git Repo & Hector Corona \\
FS-7 & TestRail Connection  & Ivonne Arizpe \\
FS-8 & Authentication & Ruby Quinones \\
FS-9 & SQL Error & Ruby Quinones \\
FS-10 & Potential Bug: User Login Issues & Hector Corona \\
\bottomrule
\end{tabular}
\end{center}
 
\section{References}
\begin{itemize}
    \item Jira Project Board: FinalProject-Scrum\\
    \url{https://calstatela-team-uqomn8ag.atlassian.net/jira/software/projects/FS/boards/133}
\end{itemize}
 
\end{document}
 


